% Volumen documental compañero. Fuente independiente, sin compilación automática.
\documentclass[a4paper,10pt]{article}
\usepackage[top=23mm,bottom=23mm,left=20mm,right=20mm,headheight=15pt]{geometry}
\usepackage{fontspec,amsmath,unicode-math}
\usepackage[spanish,es-nodecimaldot,es-noquoting]{babel}
\usepackage{microtype,booktabs,longtable,array,needspace,fancyhdr}
\usepackage[unicode,hidelinks,bookmarksopen=true,bookmarksopenlevel=1,bookmarksnumbered=false]{hyperref}
\setmainfont{STIXTwoText-Regular.otf}[BoldFont=STIXTwoText-Bold.otf,ItalicFont=STIXTwoText-Italic.otf,BoldItalicFont=STIXTwoText-BoldItalic.otf]
\setmathfont{STIXTwoMath-Regular.otf}
\renewcommand{\normalsize}{\fontsize{10.5}{13.4}\selectfont}\normalsize
\setlength{\parindent}{1em}\setlength{\parskip}{3pt}
\setlength{\emergencystretch}{2em}
\clubpenalty=10000\widowpenalty=10000\raggedbottom
\clubpenalties 3 10000 10000 0
\widowpenalties 3 10000 10000 0
\AddToHook{cmd/section/before}{\clearpage}
\AddToHook{cmd/subsection/before}{\Needspace{7\baselineskip}}
\AddToHook{cmd/subsubsection/before}{\Needspace{6\baselineskip}}
\pagestyle{fancy}\fancyhf{}
\fancyhead[L]{Catálogo estructural y metrológico de partículas}
\fancyfoot[C]{\thepage}
\hypersetup{pdftitle={Catálogo estructural y metrológico de partículas},pdfauthor={Oumar Haidara Fall; Rubén Ramos Balsa},pdfsubject={Anexo de estados, rutas y masas}}
\newcommand{\CatalogRouteReference}{Las secciones~\ref{sec:vi-324-rutas} y~\ref{sec:vi-inventario-metrologico} de este catálogo publican, respectivamente, las 324 rutas orientadas y los 855 registros metrológicos.}
\begin{document}
\thispagestyle{empty}
\begin{center}
{\fontsize{8}{10}\selectfont Manuscrito de recopilación para transferencia académica\par}
\vspace{12mm}
{\fontsize{17}{20.5}\selectfont\bfseries Catálogo estructural y metrológico de partículas\par}
\vspace{6pt}
{\fontsize{11}{13.5}\selectfont Anexo de estados, rutas y masas\par}
\vspace{11pt}
{\fontsize{11.5}{14}\selectfont Oumar Haidara Fall \textperiodcentered\ Rubén Ramos Balsa\par}
\vspace{4pt}
{\fontsize{8}{10}\selectfont Coautoría sin prelación de contribución\par}
\vspace{3pt}
{\fontsize{7.5}{9.5}\selectfont Integración y recopilación documental generadas por ChatGPT -- Astra.\par}
\end{center}
\vspace{7pt}
\begin{center}\bfseries Presentación\end{center}
\begingroup
\fontsize{10.5}{12.8}\selectfont
\setlength{\parskip}{2pt plus 1pt minus .5pt}
\noindent
\begin{minipage}{\textwidth}
Este catálogo reúne el inventario estructural de \(324\) rutas orientadas
y \(855\) registros metrológicos de masas y anchuras. Su organización
permite consultar la identidad y los campos de cada ruta y localizar las
observaciones metrológicas, sin equiparar el número de registros documentales con el
de predicciones individualizadas. Acompaña a \emph{Partículas,
persistencia y derivación del espectro de masas}, donde se desarrollan
las definiciones y demostraciones que construyen la partícula, el atlas
y el operador de masa. Aquí se publica íntegro el dominio de datos
estructurales y de contraste utilizado por ese desarrollo.

La primera sección distingue \(13\) multisecciones, \(56\) familias de
ruta y \(23\) clases del inventario. La
sección~\ref{sec:vi-324-rutas} presenta las \(324\) rutas, conservando
todos los campos de sus dos registros concordantes. La
sección~\ref{sec:vi-inventario-metrologico} reúne \(471\) masas y
\(384\) anchuras de la edición metrológica fijada, con sus identificadores,
estados, unidades y fuentes. Valores centrales, intervalos, límites y
anchuras conservan su significado propio. La presencia de una
observación en el inventario permite localizarla para el contraste;
su identificación física requiere el lector y la carta especificados
en el artículo.

Cada ficha conserva su registro íntegro, incluidos los campos de
procedencia y las distinciones entre rutas, familias y observables. El
índice enlazado, los marcadores por celda y orientación y el identificador
repetido en las páginas de continuación permiten recorrer el conjunto y
regresar directamente a la ruta elegida. La separación de este anexo de
datos mantiene legibles las fichas sin abreviar sus registros.\par
\end{minipage}
\par\endgroup
\clearpage
\tableofcontents
\clearpage
\input{sections/10a_dependencia_activa_l_g.tex}
% Generado por technical/generar_apendices_catalogo.py. No editar a mano.
\clearpage
\section{Catálogos de multisecciones, tipos de ruta y clases sectoriales}
\label{sec:vi-catalogos-estructurales}

Este apéndice organiza tres niveles distintos del inventario interno. Las
13 multisecciones reúnen celdas del atlas; las 56 colecciones clasifican rutas
internas; las 23 clases documentan la articulación con sectores y
representaciones físicas. Estos cardinales no cuentan el mismo objeto y
no se suman para obtener un número de partículas. Las tablas acompañan las
definiciones y demostraciones del artículo: su función es permitir localizar
cada registro, sus incidencias y sus lectores, no sustituir una prueba por un censo.

Se conserva el orden documental de las filas, sus identificadores y todas
las repeticiones que figuren en las fuentes. Los identificadores literales
permiten cotejar las tablas con los archivos JSON. Una raya indica una celda
vacía de la fuente; no representa un cero. Los archivos conservan además todos
los campos y cifras, los estados originales y la procedencia verificable.
\CatalogRouteReference
% La remisión depende del punto de entrada: artículo o catálogo.

\subsection{Las 13 multisecciones del atlas}

Para una multisección \(M\), sea \(\mathcal C(M)\) su conjunto de celdas y
sea \(\gamma_{B_1}(c)\) la ruta primaria que el atlas conserva en la celda
\(c\). En este apéndice estructural se define la fibra
\[
 B_1(M):=
 \operatorname{span}\bigl\{e_{\gamma_{B_1}(c)}:c\in\mathcal C(M)\bigr\}
 \subseteq \ell^2(\mathcal R_{\mathrm{adm}}).
\]
Las rutas primarias de celdas distintas determinan vectores básicos distintos;
por consiguiente,
\(\operatorname{rank}B_1(M)=\dim B_1(M)=\#\mathcal C(M)\).
Esta definición fija el objeto al que se refiere el rango de la tabla; la
columna de colecciones cuenta las
configuraciones de ruta presentes, no especies físicas. Se conservan separadamente
los espectros enteros de los lectores \(K_D\) y \(K_\Omega\).
En la notación literal \texttt{x} indica multiplicidad. Las 13 filas tienen
el estado documental de operador diagonal exacto sobre la fibra \(B_1\).
Su representación como escalar físico requiere el transporte declarado, salvo
la selección de rango uno prevista en el propio registro. Este enunciado
reproduce el alcance registrado y no convierte todos los espectros en masas únicas.

\begingroup
\setlength{\tabcolsep}{4pt}
\setlength{\extrarowheight}{4pt}
\renewcommand{\arraystretch}{1.13}
\begin{longtable}{@{}>{\raggedright\arraybackslash}p{\dimexpr 0.19871795\linewidth-2\tabcolsep\relax}>{\raggedright\arraybackslash}p{\dimexpr 0.19871795\linewidth-2\tabcolsep\relax}>{\raggedright\arraybackslash}p{\dimexpr 0.30769231\linewidth-2\tabcolsep\relax}>{\raggedright\arraybackslash}p{\dimexpr 0.29487179\linewidth-2\tabcolsep\relax}@{}}
\caption{Las 13 multisecciones, sin selección de filas.}\label{tab:vi-multisecciones}\\
\toprule
Multisección & Rango e incidencia & Clasificación y profundidad & Lectores enteros \\
\midrule
\endfirsthead
\multicolumn{4}{l}{\textit{Continuación}}\\
\toprule
Multisección & Rango e incidencia & Clasificación y profundidad & Lectores enteros \\
\midrule
\endhead
\midrule
\multicolumn{4}{r}{Continúa en la página siguiente}\\
\endfoot
\bottomrule
\endlastfoot
% VI-CATALOG-ROW multisecciones 1
charged\_\allowbreak{}lepton\_\allowbreak{}G0 & \textit{Rango}: 1\par \textit{Celdas (identificadores)}: 41\par \textit{Colecciones}: 1 & \textit{Colección}: charged\_\allowbreak{}lepton\par \textit{Profundidad}: G0\_\allowbreak{}center\par \textit{Color}: none & \textit{K D}: (80,\allowbreak{}54,\allowbreak{}6)\allowbreak{}x1\par \textit{K Ω}: (80,\allowbreak{}54,\allowbreak{}6)\allowbreak{}x1 \\
% VI-CATALOG-ROW multisecciones 2
charged\_\allowbreak{}lepton\_\allowbreak{}G2 & \textit{Rango}: 8\par \textit{Celdas (identificadores)}: 21,\allowbreak{}23,\allowbreak{}25,\allowbreak{}39,\allowbreak{}43,\allowbreak{}57,\allowbreak{}59,\allowbreak{}61\par \textit{Colecciones}: 8 & \textit{Colección}: charged\_\allowbreak{}lepton\par \textit{Profundidad}: G2\par \textit{Color}: none & \textit{K D}: (0,\allowbreak{}24,\allowbreak{}0)\allowbreak{}x1; (40,\allowbreak{}78,\allowbreak{}27)\allowbreak{}x2; (80,\allowbreak{}54,\allowbreak{}7)\allowbreak{}x1; (80,\allowbreak{}66,\allowbreak{}2)\allowbreak{}x1; (80,\allowbreak{}66,\allowbreak{}3)\allowbreak{}x1; (100,\allowbreak{}66,\allowbreak{}0)\allowbreak{}x1; (100,\allowbreak{}66,\allowbreak{}2)\allowbreak{}x1\par \textit{K Ω}: (80,\allowbreak{}54,\allowbreak{}6)\allowbreak{}x8 \\
% VI-CATALOG-ROW multisecciones 3
charged\_\allowbreak{}lepton\_\allowbreak{}G4 & \textit{Rango}: 16\par \textit{Celdas (identificadores)}: 1,\allowbreak{}3,\allowbreak{}5,\allowbreak{}7,\allowbreak{}9,\allowbreak{}19,\allowbreak{}27,\allowbreak{}37,\allowbreak{}45,\allowbreak{}55,\allowbreak{}63,\allowbreak{}73,\allowbreak{}75,\allowbreak{}77,\allowbreak{}79,\allowbreak{}81\par \textit{Colecciones}: 16 & \textit{Colección}: charged\_\allowbreak{}lepton\par \textit{Profundidad}: G4\par \textit{Color}: none & \textit{K D}: (0,\allowbreak{}24,\allowbreak{}0)\allowbreak{}x1; (40,\allowbreak{}24,\allowbreak{}2)\allowbreak{}x2; (40,\allowbreak{}54,\allowbreak{}6)\allowbreak{}x2; (40,\allowbreak{}84,\allowbreak{}30)\allowbreak{}x2; (60,\allowbreak{}78,\allowbreak{}25)\allowbreak{}x3; (80,\allowbreak{}66,\allowbreak{}2)\allowbreak{}x2; (80,\allowbreak{}66,\allowbreak{}3)\allowbreak{}x3; (80,\allowbreak{}66,\allowbreak{}4)\allowbreak{}x1\par \textit{K Ω}: (24,\allowbreak{}54,\allowbreak{}2)\allowbreak{}x1; (56,\allowbreak{}54,\allowbreak{}5)\allowbreak{}x4; (80,\allowbreak{}54,\allowbreak{}6)\allowbreak{}x11 \\
% VI-CATALOG-ROW multisecciones 4
neutrino\_\allowbreak{}G1 & \textit{Rango}: 2\par \textit{Celdas (identificadores)}: 40,\allowbreak{}42\par \textit{Colecciones}: 2 & \textit{Colección}: neutrino\par \textit{Profundidad}: G1\par \textit{Color}: none & \textit{K D}: (40,\allowbreak{}24,\allowbreak{}2)\allowbreak{}x1; (40,\allowbreak{}78,\allowbreak{}27)\allowbreak{}x1\par \textit{K Ω}: (40,\allowbreak{}54,\allowbreak{}4)\allowbreak{}x1; (80,\allowbreak{}54,\allowbreak{}6)\allowbreak{}x1 \\
% VI-CATALOG-ROW multisecciones 5
neutrino\_\allowbreak{}G2 & \textit{Rango}: 4\par \textit{Celdas (identificadores)}: 22,\allowbreak{}24,\allowbreak{}58,\allowbreak{}60\par \textit{Colecciones}: 4 & \textit{Colección}: neutrino\par \textit{Profundidad}: G2\par \textit{Color}: none & \textit{K D}: (0,\allowbreak{}24,\allowbreak{}0)\allowbreak{}x1; (40,\allowbreak{}84,\allowbreak{}30)\allowbreak{}x1; (60,\allowbreak{}78,\allowbreak{}25)\allowbreak{}x1; (80,\allowbreak{}66,\allowbreak{}3)\allowbreak{}x1\par \textit{K Ω}: (4,\allowbreak{}54,\allowbreak{}2)\allowbreak{}x1; (80,\allowbreak{}54,\allowbreak{}6)\allowbreak{}x3 \\
% VI-CATALOG-ROW multisecciones 6
neutrino\_\allowbreak{}G3 & \textit{Rango}: 6\par \textit{Celdas (identificadores)}: 20,\allowbreak{}26,\allowbreak{}38,\allowbreak{}44,\allowbreak{}56,\allowbreak{}62\par \textit{Colecciones}: 6 & \textit{Colección}: neutrino\par \textit{Profundidad}: G3\par \textit{Color}: none & \textit{K D}: (40,\allowbreak{}24,\allowbreak{}2)\allowbreak{}x2; (40,\allowbreak{}78,\allowbreak{}27)\allowbreak{}x1; (60,\allowbreak{}84,\allowbreak{}28)\allowbreak{}x1; (80,\allowbreak{}54,\allowbreak{}6)\allowbreak{}x1; (80,\allowbreak{}66,\allowbreak{}3)\allowbreak{}x1\par \textit{K Ω}: (36,\allowbreak{}54,\allowbreak{}4)\allowbreak{}x1; (56,\allowbreak{}54,\allowbreak{}5)\allowbreak{}x1; (80,\allowbreak{}54,\allowbreak{}6)\allowbreak{}x4 \\
% VI-CATALOG-ROW multisecciones 7
neutrino\_\allowbreak{}G4 & \textit{Rango}: 8\par \textit{Celdas (identificadores)}: 2,\allowbreak{}4,\allowbreak{}6,\allowbreak{}8,\allowbreak{}74,\allowbreak{}76,\allowbreak{}78,\allowbreak{}80\par \textit{Colecciones}: 8 & \textit{Colección}: neutrino\par \textit{Profundidad}: G4\par \textit{Color}: none & \textit{K D}: (0,\allowbreak{}24,\allowbreak{}0)\allowbreak{}x1; (40,\allowbreak{}24,\allowbreak{}2)\allowbreak{}x1; (40,\allowbreak{}78,\allowbreak{}27)\allowbreak{}x2; (40,\allowbreak{}84,\allowbreak{}30)\allowbreak{}x1; (80,\allowbreak{}54,\allowbreak{}7)\allowbreak{}x1; (80,\allowbreak{}66,\allowbreak{}2)\allowbreak{}x1; (100,\allowbreak{}66,\allowbreak{}0)\allowbreak{}x1\par \textit{K Ω}: (36,\allowbreak{}54,\allowbreak{}4)\allowbreak{}x1; (56,\allowbreak{}54,\allowbreak{}5)\allowbreak{}x1; (80,\allowbreak{}54,\allowbreak{}6)\allowbreak{}x6 \\
% VI-CATALOG-ROW multisecciones 8
quark\_\allowbreak{}down\_\allowbreak{}G1 & \textit{Rango}: 2\par \textit{Celdas (identificadores)}: 32,\allowbreak{}50\par \textit{Colecciones}: 2 & \textit{Colección}: quark\_\allowbreak{}down\par \textit{Profundidad}: G1\par \textit{Color}: B | G & \textit{K D}: (40,\allowbreak{}24,\allowbreak{}2)\allowbreak{}x1; (40,\allowbreak{}78,\allowbreak{}27)\allowbreak{}x1\par \textit{K Ω}: (40,\allowbreak{}54,\allowbreak{}4)\allowbreak{}x1; (80,\allowbreak{}54,\allowbreak{}6)\allowbreak{}x1 \\
% VI-CATALOG-ROW multisecciones 9
quark\_\allowbreak{}down\_\allowbreak{}G2 & \textit{Rango}: 4\par \textit{Celdas (identificadores)}: 30,\allowbreak{}34,\allowbreak{}48,\allowbreak{}52\par \textit{Colecciones}: 4 & \textit{Colección}: quark\_\allowbreak{}down\par \textit{Profundidad}: G2\par \textit{Color}: B | G | R & \textit{K D}: (0,\allowbreak{}24,\allowbreak{}0)\allowbreak{}x1; (40,\allowbreak{}84,\allowbreak{}30)\allowbreak{}x1; (60,\allowbreak{}78,\allowbreak{}25)\allowbreak{}x1; (80,\allowbreak{}66,\allowbreak{}2)\allowbreak{}x1\par \textit{K Ω}: (4,\allowbreak{}54,\allowbreak{}2)\allowbreak{}x1; (80,\allowbreak{}54,\allowbreak{}6)\allowbreak{}x3 \\
% VI-CATALOG-ROW multisecciones 10
quark\_\allowbreak{}down\_\allowbreak{}G3 & \textit{Rango}: 6\par \textit{Celdas (identificadores)}: 12,\allowbreak{}14,\allowbreak{}16,\allowbreak{}66,\allowbreak{}68,\allowbreak{}70\par \textit{Colecciones}: 6 & \textit{Colección}: quark\_\allowbreak{}down\par \textit{Profundidad}: G3\par \textit{Color}: B | G | R & \textit{K D}: (40,\allowbreak{}24,\allowbreak{}2)\allowbreak{}x2; (40,\allowbreak{}78,\allowbreak{}27)\allowbreak{}x1; (60,\allowbreak{}84,\allowbreak{}28)\allowbreak{}x1; (80,\allowbreak{}54,\allowbreak{}6)\allowbreak{}x1; (80,\allowbreak{}66,\allowbreak{}4)\allowbreak{}x1\par \textit{K Ω}: (36,\allowbreak{}54,\allowbreak{}4)\allowbreak{}x1; (80,\allowbreak{}54,\allowbreak{}6)\allowbreak{}x5 \\
% VI-CATALOG-ROW multisecciones 11
quark\_\allowbreak{}down\_\allowbreak{}G4 & \textit{Rango}: 8\par \textit{Celdas (identificadores)}: 10,\allowbreak{}18,\allowbreak{}28,\allowbreak{}36,\allowbreak{}46,\allowbreak{}54,\allowbreak{}64,\allowbreak{}72\par \textit{Colecciones}: 8 & \textit{Colección}: quark\_\allowbreak{}down\par \textit{Profundidad}: G4\par \textit{Color}: B | G | R & \textit{K D}: (0,\allowbreak{}24,\allowbreak{}0)\allowbreak{}x1; (40,\allowbreak{}24,\allowbreak{}2)\allowbreak{}x1; (40,\allowbreak{}78,\allowbreak{}27)\allowbreak{}x2; (40,\allowbreak{}84,\allowbreak{}30)\allowbreak{}x1; (80,\allowbreak{}54,\allowbreak{}7)\allowbreak{}x1; (80,\allowbreak{}66,\allowbreak{}3)\allowbreak{}x1; (100,\allowbreak{}66,\allowbreak{}2)\allowbreak{}x1\par \textit{K Ω}: (36,\allowbreak{}54,\allowbreak{}4)\allowbreak{}x1; (80,\allowbreak{}54,\allowbreak{}6)\allowbreak{}x7 \\
% VI-CATALOG-ROW multisecciones 12
quark\_\allowbreak{}up\_\allowbreak{}G1 & \textit{Rango}: 4\par \textit{Celdas (identificadores)}: 31,\allowbreak{}33,\allowbreak{}49,\allowbreak{}51\par \textit{Colecciones}: 3 & \textit{Colección}: quark\_\allowbreak{}up\par \textit{Profundidad}: G1\par \textit{Color}: B | G | R & \textit{K D}: (40,\allowbreak{}54,\allowbreak{}6)\allowbreak{}x1; (60,\allowbreak{}78,\allowbreak{}25)\allowbreak{}x1; (80,\allowbreak{}66,\allowbreak{}2)\allowbreak{}x1; (80,\allowbreak{}66,\allowbreak{}3)\allowbreak{}x1\par \textit{K Ω}: (80,\allowbreak{}54,\allowbreak{}6)\allowbreak{}x4 \\
% VI-CATALOG-ROW multisecciones 13
quark\_\allowbreak{}up\_\allowbreak{}G3 & \textit{Rango}: 12\par \textit{Celdas (identificadores)}: 11,\allowbreak{}13,\allowbreak{}15,\allowbreak{}17,\allowbreak{}29,\allowbreak{}35,\allowbreak{}47,\allowbreak{}53,\allowbreak{}65,\allowbreak{}67,\allowbreak{}69,\allowbreak{}71\par \textit{Colecciones}: 12 & \textit{Colección}: quark\_\allowbreak{}up\par \textit{Profundidad}: G3\par \textit{Color}: B | G | R & \textit{K D}: (40,\allowbreak{}24,\allowbreak{}2)\allowbreak{}x2; (40,\allowbreak{}78,\allowbreak{}27)\allowbreak{}x2; (60,\allowbreak{}84,\allowbreak{}28)\allowbreak{}x1; (80,\allowbreak{}54,\allowbreak{}6)\allowbreak{}x3; (80,\allowbreak{}66,\allowbreak{}3)\allowbreak{}x1; (80,\allowbreak{}66,\allowbreak{}4)\allowbreak{}x1; (100,\allowbreak{}66,\allowbreak{}0)\allowbreak{}x1; (100,\allowbreak{}66,\allowbreak{}2)\allowbreak{}x1\par \textit{K Ω}: (56,\allowbreak{}54,\allowbreak{}5)\allowbreak{}x3; (80,\allowbreak{}54,\allowbreak{}6)\allowbreak{}x9 \\
\end{longtable}
\endgroup
\subsection{Clasificación de los 56 tipos de ruta}

La firma que identifica una colección se conserva literalmente, con sus signos
y separadores. Las celdas y multisecciones hacen visible que una colección de
ruta no equivale por definición a una especie física. Los perfiles conjuntos
cuentan las combinaciones registradas dentro de la colección. Los espectros
decimales del carácter y los espectros completos de firma permanecen íntegros
en \texttt{catalogo\_familias.json}; aquí se expresan la incidencia y los
lectores enteros que permiten orientarse en esa información.

\begingroup
\setlength{\tabcolsep}{4pt}
\setlength{\extrarowheight}{4pt}
\renewcommand{\arraystretch}{1.13}
\begin{longtable}{@{}>{\raggedright\arraybackslash}p{\dimexpr 0.21794872\linewidth-2\tabcolsep\relax}>{\raggedright\arraybackslash}p{\dimexpr 0.34615385\linewidth-2\tabcolsep\relax}>{\raggedright\arraybackslash}p{\dimexpr 0.21794872\linewidth-2\tabcolsep\relax}>{\raggedright\arraybackslash}p{\dimexpr 0.21794872\linewidth-2\tabcolsep\relax}@{}}
\caption{Las 56 colecciones internas de ruta.}\label{tab:vi-familias}\\
\toprule
Colección y firma & Incidencia & \(K_D\) & \(K_\Omega\) \\
\midrule
\endfirsthead
\multicolumn{4}{l}{\textit{Continuación}}\\
\toprule
Colección y firma & Incidencia & \(K_D\) & \(K_\Omega\) \\
\midrule
\endhead
\midrule
\multicolumn{4}{r}{Continúa en la página siguiente}\\
\endfoot
\bottomrule
\endlastfoot
% VI-CATALOG-ROW familias 1
1\par 24|\allowbreak{}0|\allowbreak{}12|\allowbreak{}0|\allowbreak{}-12|\allowbreak{}+++-\kern0pt{}-\kern0pt{}-+++-\kern0pt{}-\kern0pt{}- & \textit{Número de celdas}: 5\par \textit{Celdas (identificadores)}: 21,\allowbreak{}31,\allowbreak{}41,\allowbreak{}51,\allowbreak{}81\par \textit{Multisecciones}: charged\_\allowbreak{}lepton\_\allowbreak{}G0 | charged\_\allowbreak{}lepton\_\allowbreak{}G2 | charged\_\allowbreak{}lepton\_\allowbreak{}G4 | quark\_\allowbreak{}up\_\allowbreak{}G1\par \textit{Colecciones}: charged\_\allowbreak{}lepton | quark\_\allowbreak{}up\par \textit{Color}: R | none\par \textit{Perfiles conjuntos}: 5 & (0,\allowbreak{}24,\allowbreak{}0)\allowbreak{}x2; (40,\allowbreak{}54,\allowbreak{}6)\allowbreak{}x1; (60,\allowbreak{}78,\allowbreak{}25)\allowbreak{}x1; (80,\allowbreak{}54,\allowbreak{}6)\allowbreak{}x1 & (24,\allowbreak{}54,\allowbreak{}2)\allowbreak{}x1; (80,\allowbreak{}54,\allowbreak{}6)\allowbreak{}x4 \\
% VI-CATALOG-ROW familias 2
2\par 24|\allowbreak{}0|\allowbreak{}4|\allowbreak{}0|\allowbreak{}-12|\allowbreak{}+++-\kern0pt{}-\kern0pt{}-+++-\kern0pt{}-\kern0pt{}- & \textit{Número de celdas}: 2\par \textit{Celdas (identificadores)}: 11,\allowbreak{}61\par \textit{Multisecciones}: charged\_\allowbreak{}lepton\_\allowbreak{}G2 | quark\_\allowbreak{}up\_\allowbreak{}G3\par \textit{Colecciones}: charged\_\allowbreak{}lepton | quark\_\allowbreak{}up\par \textit{Color}: R | none\par \textit{Perfiles conjuntos}: 2 & (60,\allowbreak{}84,\allowbreak{}28)\allowbreak{}x1; (80,\allowbreak{}54,\allowbreak{}7)\allowbreak{}x1 & (80,\allowbreak{}54,\allowbreak{}6)\allowbreak{}x2 \\
% VI-CATALOG-ROW familias 3
3\par 24|\allowbreak{}0|\allowbreak{}8|\allowbreak{}0|\allowbreak{}-12|\allowbreak{}+++-\kern0pt{}-\kern0pt{}-+++-\kern0pt{}-\kern0pt{}- & \textit{Número de celdas}: 2\par \textit{Celdas (identificadores)}: 1,\allowbreak{}71\par \textit{Multisecciones}: charged\_\allowbreak{}lepton\_\allowbreak{}G4 | quark\_\allowbreak{}up\_\allowbreak{}G3\par \textit{Colecciones}: charged\_\allowbreak{}lepton | quark\_\allowbreak{}up\par \textit{Color}: R | none\par \textit{Perfiles conjuntos}: 2 & (60,\allowbreak{}78,\allowbreak{}25)\allowbreak{}x1; (80,\allowbreak{}54,\allowbreak{}6)\allowbreak{}x1 & (80,\allowbreak{}54,\allowbreak{}6)\allowbreak{}x2 \\
% VI-CATALOG-ROW familias 4
4\par 28|\allowbreak{}-6|\allowbreak{}-2|\allowbreak{}-2|\allowbreak{}-8|\allowbreak{}0++-\kern0pt{}-\kern0pt{}-0++-\kern0pt{}-\kern0pt{}- & \textit{Número de celdas}: 1\par \textit{Celdas (identificadores)}: 53\par \textit{Multisecciones}: quark\_\allowbreak{}up\_\allowbreak{}G3\par \textit{Colecciones}: quark\_\allowbreak{}up\par \textit{Color}: G\par \textit{Perfiles conjuntos}: 1 & (80,\allowbreak{}66,\allowbreak{}3)\allowbreak{}x1 & (56,\allowbreak{}54,\allowbreak{}5)\allowbreak{}x1 \\
% VI-CATALOG-ROW familias 5
5\par 28|\allowbreak{}-6|\allowbreak{}-4|\allowbreak{}-2|\allowbreak{}-8|\allowbreak{}0++-\kern0pt{}-\kern0pt{}-0++-\kern0pt{}-\kern0pt{}- & \textit{Número de celdas}: 1\par \textit{Celdas (identificadores)}: 20\par \textit{Multisecciones}: neutrino\_\allowbreak{}G3\par \textit{Colecciones}: neutrino\par \textit{Color}: none\par \textit{Perfiles conjuntos}: 1 & (80,\allowbreak{}66,\allowbreak{}3)\allowbreak{}x1 & (56,\allowbreak{}54,\allowbreak{}5)\allowbreak{}x1 \\
% VI-CATALOG-ROW familias 6
6\par 28|\allowbreak{}-6|\allowbreak{}0|\allowbreak{}0|\allowbreak{}-12|\allowbreak{}+++-\kern0pt{}-\kern0pt{}-+++-\kern0pt{}-\kern0pt{}- & \textit{Número de celdas}: 2\par \textit{Celdas (identificadores)}: 30,\allowbreak{}70\par \textit{Multisecciones}: quark\_\allowbreak{}down\_\allowbreak{}G2 | quark\_\allowbreak{}down\_\allowbreak{}G3\par \textit{Colecciones}: quark\_\allowbreak{}down\par \textit{Color}: G\par \textit{Perfiles conjuntos}: 2 & (40,\allowbreak{}24,\allowbreak{}2)\allowbreak{}x1; (80,\allowbreak{}66,\allowbreak{}2)\allowbreak{}x1 & (80,\allowbreak{}54,\allowbreak{}6)\allowbreak{}x2 \\
% VI-CATALOG-ROW familias 7
7\par 28|\allowbreak{}-6|\allowbreak{}10|\allowbreak{}0|\allowbreak{}-12|\allowbreak{}+++-\kern0pt{}-\kern0pt{}-+++-\kern0pt{}-\kern0pt{}- & \textit{Número de celdas}: 2\par \textit{Celdas (identificadores)}: 33,\allowbreak{}64\par \textit{Multisecciones}: quark\_\allowbreak{}down\_\allowbreak{}G4 | quark\_\allowbreak{}up\_\allowbreak{}G1\par \textit{Colecciones}: quark\_\allowbreak{}down | quark\_\allowbreak{}up\par \textit{Color}: G\par \textit{Perfiles conjuntos}: 2 & (40,\allowbreak{}24,\allowbreak{}2)\allowbreak{}x1; (80,\allowbreak{}66,\allowbreak{}2)\allowbreak{}x1 & (36,\allowbreak{}54,\allowbreak{}4)\allowbreak{}x1; (80,\allowbreak{}54,\allowbreak{}6)\allowbreak{}x1 \\
% VI-CATALOG-ROW familias 8
8\par 28|\allowbreak{}-6|\allowbreak{}12|\allowbreak{}0|\allowbreak{}-12|\allowbreak{}+++-\kern0pt{}-\kern0pt{}-+++-\kern0pt{}-\kern0pt{}- & \textit{Número de celdas}: 1\par \textit{Celdas (identificadores)}: 10\par \textit{Multisecciones}: quark\_\allowbreak{}down\_\allowbreak{}G4\par \textit{Colecciones}: quark\_\allowbreak{}down\par \textit{Color}: G\par \textit{Perfiles conjuntos}: 1 & (100,\allowbreak{}66,\allowbreak{}2)\allowbreak{}x1 & (80,\allowbreak{}54,\allowbreak{}6)\allowbreak{}x1 \\
% VI-CATALOG-ROW familias 9
9\par 28|\allowbreak{}-6|\allowbreak{}2|\allowbreak{}0|\allowbreak{}-12|\allowbreak{}+++-\kern0pt{}-\kern0pt{}-+++-\kern0pt{}-\kern0pt{}- & \textit{Número de celdas}: 1\par \textit{Celdas (identificadores)}: 3\par \textit{Multisecciones}: charged\_\allowbreak{}lepton\_\allowbreak{}G4\par \textit{Colecciones}: charged\_\allowbreak{}lepton\par \textit{Color}: none\par \textit{Perfiles conjuntos}: 1 & (40,\allowbreak{}84,\allowbreak{}30)\allowbreak{}x1 & (80,\allowbreak{}54,\allowbreak{}6)\allowbreak{}x1 \\
% VI-CATALOG-ROW familias 10
10\par 28|\allowbreak{}-6|\allowbreak{}4|\allowbreak{}-2|\allowbreak{}-8|\allowbreak{}++0-\kern0pt{}-\kern0pt{}-++0-\kern0pt{}-\kern0pt{}- & \textit{Número de celdas}: 1\par \textit{Celdas (identificadores)}: 50\par \textit{Multisecciones}: quark\_\allowbreak{}down\_\allowbreak{}G1\par \textit{Colecciones}: quark\_\allowbreak{}down\par \textit{Color}: G\par \textit{Perfiles conjuntos}: 1 & (40,\allowbreak{}78,\allowbreak{}27)\allowbreak{}x1 & (80,\allowbreak{}54,\allowbreak{}6)\allowbreak{}x1 \\
% VI-CATALOG-ROW familias 11
11\par 28|\allowbreak{}-6|\allowbreak{}4|\allowbreak{}0|\allowbreak{}-12|\allowbreak{}+++-\kern0pt{}-\kern0pt{}-+++-\kern0pt{}-\kern0pt{}- & \textit{Número de celdas}: 2\par \textit{Celdas (identificadores)}: 9,\allowbreak{}40\par \textit{Multisecciones}: charged\_\allowbreak{}lepton\_\allowbreak{}G4 | neutrino\_\allowbreak{}G1\par \textit{Colecciones}: charged\_\allowbreak{}lepton | neutrino\par \textit{Color}: none\par \textit{Perfiles conjuntos}: 2 & (40,\allowbreak{}24,\allowbreak{}2)\allowbreak{}x1; (80,\allowbreak{}66,\allowbreak{}2)\allowbreak{}x1 & (40,\allowbreak{}54,\allowbreak{}4)\allowbreak{}x1; (80,\allowbreak{}54,\allowbreak{}6)\allowbreak{}x1 \\
% VI-CATALOG-ROW familias 12
12\par 28|\allowbreak{}-6|\allowbreak{}4|\allowbreak{}0|\allowbreak{}-8|\allowbreak{}+0+-\kern0pt{}-\kern0pt{}-+0+-\kern0pt{}-\kern0pt{}- & \textit{Número de celdas}: 1\par \textit{Celdas (identificadores)}: 80\par \textit{Multisecciones}: neutrino\_\allowbreak{}G4\par \textit{Colecciones}: neutrino\par \textit{Color}: none\par \textit{Perfiles conjuntos}: 1 & (40,\allowbreak{}78,\allowbreak{}27)\allowbreak{}x1 & (80,\allowbreak{}54,\allowbreak{}6)\allowbreak{}x1 \\
% VI-CATALOG-ROW familias 13
13\par 28|\allowbreak{}-6|\allowbreak{}6|\allowbreak{}-2|\allowbreak{}-8|\allowbreak{}++0-\kern0pt{}-\kern0pt{}-++0-\kern0pt{}-\kern0pt{}- & \textit{Número de celdas}: 1\par \textit{Celdas (identificadores)}: 74\par \textit{Multisecciones}: neutrino\_\allowbreak{}G4\par \textit{Colecciones}: neutrino\par \textit{Color}: none\par \textit{Perfiles conjuntos}: 1 & (40,\allowbreak{}78,\allowbreak{}27)\allowbreak{}x1 & (80,\allowbreak{}54,\allowbreak{}6)\allowbreak{}x1 \\
% VI-CATALOG-ROW familias 14
14\par 28|\allowbreak{}-6|\allowbreak{}6|\allowbreak{}0|\allowbreak{}-12|\allowbreak{}+++-\kern0pt{}-\kern0pt{}-+++-\kern0pt{}-\kern0pt{}- & \textit{Número de celdas}: 3\par \textit{Celdas (identificadores)}: 13,\allowbreak{}43,\allowbreak{}63\par \textit{Multisecciones}: charged\_\allowbreak{}lepton\_\allowbreak{}G2 | charged\_\allowbreak{}lepton\_\allowbreak{}G4 | quark\_\allowbreak{}up\_\allowbreak{}G3\par \textit{Colecciones}: charged\_\allowbreak{}lepton | quark\_\allowbreak{}up\par \textit{Color}: G | none\par \textit{Perfiles conjuntos}: 3 & (40,\allowbreak{}24,\allowbreak{}2)\allowbreak{}x1; (80,\allowbreak{}66,\allowbreak{}2)\allowbreak{}x1; (100,\allowbreak{}66,\allowbreak{}2)\allowbreak{}x1 & (80,\allowbreak{}54,\allowbreak{}6)\allowbreak{}x3 \\
% VI-CATALOG-ROW familias 15
15\par 28|\allowbreak{}-6|\allowbreak{}6|\allowbreak{}0|\allowbreak{}-8|\allowbreak{}+0+-\kern0pt{}-\kern0pt{}-+0+-\kern0pt{}-\kern0pt{}- & \textit{Número de celdas}: 1\par \textit{Celdas (identificadores)}: 23\par \textit{Multisecciones}: charged\_\allowbreak{}lepton\_\allowbreak{}G2\par \textit{Colecciones}: charged\_\allowbreak{}lepton\par \textit{Color}: none\par \textit{Perfiles conjuntos}: 1 & (40,\allowbreak{}78,\allowbreak{}27)\allowbreak{}x1 & (80,\allowbreak{}54,\allowbreak{}6)\allowbreak{}x1 \\
% VI-CATALOG-ROW familias 16
16\par 28|\allowbreak{}-6|\allowbreak{}8|\allowbreak{}0|\allowbreak{}-12|\allowbreak{}+++-\kern0pt{}-\kern0pt{}-+++-\kern0pt{}-\kern0pt{}- & \textit{Número de celdas}: 1\par \textit{Celdas (identificadores)}: 60\par \textit{Multisecciones}: neutrino\_\allowbreak{}G2\par \textit{Colecciones}: neutrino\par \textit{Color}: none\par \textit{Perfiles conjuntos}: 1 & (40,\allowbreak{}84,\allowbreak{}30)\allowbreak{}x1 & (80,\allowbreak{}54,\allowbreak{}6)\allowbreak{}x1 \\
% VI-CATALOG-ROW familias 17
17\par 28|\allowbreak{}6|\allowbreak{}-10|\allowbreak{}0|\allowbreak{}-12|\allowbreak{}+++-\kern0pt{}-\kern0pt{}-+++-\kern0pt{}-\kern0pt{}- & \textit{Número de celdas}: 1\par \textit{Celdas (identificadores)}: 19\par \textit{Multisecciones}: charged\_\allowbreak{}lepton\_\allowbreak{}G4\par \textit{Colecciones}: charged\_\allowbreak{}lepton\par \textit{Color}: none\par \textit{Perfiles conjuntos}: 1 & (40,\allowbreak{}84,\allowbreak{}30)\allowbreak{}x1 & (80,\allowbreak{}54,\allowbreak{}6)\allowbreak{}x1 \\
% VI-CATALOG-ROW familias 18
18\par 28|\allowbreak{}6|\allowbreak{}-2|\allowbreak{}0|\allowbreak{}-12|\allowbreak{}+++-\kern0pt{}-\kern0pt{}-+++-\kern0pt{}-\kern0pt{}- & \textit{Número de celdas}: 5\par \textit{Celdas (identificadores)}: 18,\allowbreak{}29,\allowbreak{}49,\allowbreak{}59,\allowbreak{}79\par \textit{Multisecciones}: charged\_\allowbreak{}lepton\_\allowbreak{}G2 | charged\_\allowbreak{}lepton\_\allowbreak{}G4 | quark\_\allowbreak{}down\_\allowbreak{}G4 | quark\_\allowbreak{}up\_\allowbreak{}G1 | quark\_\allowbreak{}up\_\allowbreak{}G3\par \textit{Colecciones}: charged\_\allowbreak{}lepton | quark\_\allowbreak{}down | quark\_\allowbreak{}up\par \textit{Color}: B | none\par \textit{Perfiles conjuntos}: 4 & (40,\allowbreak{}24,\allowbreak{}2)\allowbreak{}x1; (40,\allowbreak{}78,\allowbreak{}27)\allowbreak{}x1; (80,\allowbreak{}66,\allowbreak{}3)\allowbreak{}x2; (100,\allowbreak{}66,\allowbreak{}0)\allowbreak{}x1 & (80,\allowbreak{}54,\allowbreak{}6)\allowbreak{}x5 \\
% VI-CATALOG-ROW familias 19
19\par 28|\allowbreak{}6|\allowbreak{}-4|\allowbreak{}0|\allowbreak{}-12|\allowbreak{}+++-\kern0pt{}-\kern0pt{}-+++-\kern0pt{}-\kern0pt{}- & \textit{Número de celdas}: 2\par \textit{Celdas (identificadores)}: 2,\allowbreak{}42\par \textit{Multisecciones}: neutrino\_\allowbreak{}G1 | neutrino\_\allowbreak{}G4\par \textit{Colecciones}: neutrino\par \textit{Color}: none\par \textit{Perfiles conjuntos}: 2 & (40,\allowbreak{}78,\allowbreak{}27)\allowbreak{}x1; (100,\allowbreak{}66,\allowbreak{}0)\allowbreak{}x1 & (80,\allowbreak{}54,\allowbreak{}6)\allowbreak{}x2 \\
% VI-CATALOG-ROW familias 20
20\par 28|\allowbreak{}6|\allowbreak{}-6|\allowbreak{}0|\allowbreak{}-12|\allowbreak{}+++-\kern0pt{}-\kern0pt{}-+++-\kern0pt{}-\kern0pt{}- & \textit{Número de celdas}: 1\par \textit{Celdas (identificadores)}: 69\par \textit{Multisecciones}: quark\_\allowbreak{}up\_\allowbreak{}G3\par \textit{Colecciones}: quark\_\allowbreak{}up\par \textit{Color}: B\par \textit{Perfiles conjuntos}: 1 & (80,\allowbreak{}66,\allowbreak{}4)\allowbreak{}x1 & (80,\allowbreak{}54,\allowbreak{}6)\allowbreak{}x1 \\
% VI-CATALOG-ROW familias 21
21\par 28|\allowbreak{}6|\allowbreak{}-8|\allowbreak{}0|\allowbreak{}-12|\allowbreak{}+++-\kern0pt{}-\kern0pt{}-+++-\kern0pt{}-\kern0pt{}- & \textit{Número de celdas}: 1\par \textit{Celdas (identificadores)}: 52\par \textit{Multisecciones}: quark\_\allowbreak{}down\_\allowbreak{}G2\par \textit{Colecciones}: quark\_\allowbreak{}down\par \textit{Color}: B\par \textit{Perfiles conjuntos}: 1 & (40,\allowbreak{}84,\allowbreak{}30)\allowbreak{}x1 & (80,\allowbreak{}54,\allowbreak{}6)\allowbreak{}x1 \\
% VI-CATALOG-ROW familias 22
22\par 28|\allowbreak{}6|\allowbreak{}0|\allowbreak{}0|\allowbreak{}-12|\allowbreak{}+++-\kern0pt{}-\kern0pt{}-+++-\kern0pt{}-\kern0pt{}- & \textit{Número de celdas}: 4\par \textit{Celdas (identificadores)}: 22,\allowbreak{}32,\allowbreak{}72,\allowbreak{}73\par \textit{Multisecciones}: charged\_\allowbreak{}lepton\_\allowbreak{}G4 | neutrino\_\allowbreak{}G2 | quark\_\allowbreak{}down\_\allowbreak{}G1 | quark\_\allowbreak{}down\_\allowbreak{}G4\par \textit{Colecciones}: charged\_\allowbreak{}lepton | neutrino | quark\_\allowbreak{}down\par \textit{Color}: B | none\par \textit{Perfiles conjuntos}: 3 & (40,\allowbreak{}24,\allowbreak{}2)\allowbreak{}x1; (40,\allowbreak{}78,\allowbreak{}27)\allowbreak{}x1; (80,\allowbreak{}66,\allowbreak{}3)\allowbreak{}x2 & (40,\allowbreak{}54,\allowbreak{}4)\allowbreak{}x1; (80,\allowbreak{}54,\allowbreak{}6)\allowbreak{}x3 \\
% VI-CATALOG-ROW familias 23
23\par 28|\allowbreak{}6|\allowbreak{}2|\allowbreak{}0|\allowbreak{}-12|\allowbreak{}+++-\kern0pt{}-\kern0pt{}-+++-\kern0pt{}-\kern0pt{}- & \textit{Número de celdas}: 2\par \textit{Celdas (identificadores)}: 8,\allowbreak{}39\par \textit{Multisecciones}: charged\_\allowbreak{}lepton\_\allowbreak{}G2 | neutrino\_\allowbreak{}G4\par \textit{Colecciones}: charged\_\allowbreak{}lepton | neutrino\par \textit{Color}: none\par \textit{Perfiles conjuntos}: 2 & (40,\allowbreak{}24,\allowbreak{}2)\allowbreak{}x1; (40,\allowbreak{}78,\allowbreak{}27)\allowbreak{}x1 & (36,\allowbreak{}54,\allowbreak{}4)\allowbreak{}x1; (80,\allowbreak{}54,\allowbreak{}6)\allowbreak{}x1 \\
% VI-CATALOG-ROW familias 24
24\par 28|\allowbreak{}6|\allowbreak{}4|\allowbreak{}0|\allowbreak{}-12|\allowbreak{}+++-\kern0pt{}-\kern0pt{}-+++-\kern0pt{}-\kern0pt{}- & \textit{Número de celdas}: 1\par \textit{Celdas (identificadores)}: 12\par \textit{Multisecciones}: quark\_\allowbreak{}down\_\allowbreak{}G3\par \textit{Colecciones}: quark\_\allowbreak{}down\par \textit{Color}: B\par \textit{Perfiles conjuntos}: 1 & (80,\allowbreak{}66,\allowbreak{}4)\allowbreak{}x1 & (80,\allowbreak{}54,\allowbreak{}6)\allowbreak{}x1 \\
% VI-CATALOG-ROW familias 25
25\par 28|\allowbreak{}6|\allowbreak{}8|\allowbreak{}0|\allowbreak{}-12|\allowbreak{}+++-\kern0pt{}-\kern0pt{}-+++-\kern0pt{}-\kern0pt{}- & \textit{Número de celdas}: 1\par \textit{Celdas (identificadores)}: 62\par \textit{Multisecciones}: neutrino\_\allowbreak{}G3\par \textit{Colecciones}: neutrino\par \textit{Color}: none\par \textit{Perfiles conjuntos}: 1 & (40,\allowbreak{}24,\allowbreak{}2)\allowbreak{}x1 & (80,\allowbreak{}54,\allowbreak{}6)\allowbreak{}x1 \\
% VI-CATALOG-ROW familias 26
26\par 30|\allowbreak{}-6|\allowbreak{}-4|\allowbreak{}-2|\allowbreak{}-8|\allowbreak{}0++-\kern0pt{}-\kern0pt{}-0++-\kern0pt{}-\kern0pt{}- & \textit{Número de celdas}: 1\par \textit{Celdas (identificadores)}: 4\par \textit{Multisecciones}: neutrino\_\allowbreak{}G4\par \textit{Colecciones}: neutrino\par \textit{Color}: none\par \textit{Perfiles conjuntos}: 1 & (80,\allowbreak{}54,\allowbreak{}7)\allowbreak{}x1 & (56,\allowbreak{}54,\allowbreak{}5)\allowbreak{}x1 \\
% VI-CATALOG-ROW familias 27
27\par 30|\allowbreak{}-6|\allowbreak{}-8|\allowbreak{}0|\allowbreak{}-8|\allowbreak{}+0+-\kern0pt{}-\kern0pt{}-+0+-\kern0pt{}-\kern0pt{}- & \textit{Número de celdas}: 1\par \textit{Celdas (identificadores)}: 55\par \textit{Multisecciones}: charged\_\allowbreak{}lepton\_\allowbreak{}G4\par \textit{Colecciones}: charged\_\allowbreak{}lepton\par \textit{Color}: none\par \textit{Perfiles conjuntos}: 1 & (60,\allowbreak{}78,\allowbreak{}25)\allowbreak{}x1 & (80,\allowbreak{}54,\allowbreak{}6)\allowbreak{}x1 \\
% VI-CATALOG-ROW familias 28
28\par 30|\allowbreak{}-6|\allowbreak{}0|\allowbreak{}-2|\allowbreak{}-8|\allowbreak{}++0-\kern0pt{}-\kern0pt{}-++0-\kern0pt{}-\kern0pt{}- & \textit{Número de celdas}: 2\par \textit{Celdas (identificadores)}: 24,\allowbreak{}34\par \textit{Multisecciones}: neutrino\_\allowbreak{}G2 | quark\_\allowbreak{}down\_\allowbreak{}G2\par \textit{Colecciones}: neutrino | quark\_\allowbreak{}down\par \textit{Color}: R | none\par \textit{Perfiles conjuntos}: 2 & (0,\allowbreak{}24,\allowbreak{}0)\allowbreak{}x1; (60,\allowbreak{}78,\allowbreak{}25)\allowbreak{}x1 & (4,\allowbreak{}54,\allowbreak{}2)\allowbreak{}x1; (80,\allowbreak{}54,\allowbreak{}6)\allowbreak{}x1 \\
% VI-CATALOG-ROW familias 29
29\par 30|\allowbreak{}-6|\allowbreak{}0|\allowbreak{}-2|\allowbreak{}-8|\allowbreak{}0++-\kern0pt{}-\kern0pt{}-0++-\kern0pt{}-\kern0pt{}- & \textit{Número de celdas}: 1\par \textit{Celdas (identificadores)}: 75\par \textit{Multisecciones}: charged\_\allowbreak{}lepton\_\allowbreak{}G4\par \textit{Colecciones}: charged\_\allowbreak{}lepton\par \textit{Color}: none\par \textit{Perfiles conjuntos}: 1 & (40,\allowbreak{}54,\allowbreak{}6)\allowbreak{}x1 & (56,\allowbreak{}54,\allowbreak{}5)\allowbreak{}x1 \\
% VI-CATALOG-ROW familias 30
30\par 30|\allowbreak{}-6|\allowbreak{}0|\allowbreak{}0|\allowbreak{}-12|\allowbreak{}+++-\kern0pt{}-\kern0pt{}-+++-\kern0pt{}-\kern0pt{}- & \textit{Número de celdas}: 1\par \textit{Celdas (identificadores)}: 14\par \textit{Multisecciones}: quark\_\allowbreak{}down\_\allowbreak{}G3\par \textit{Colecciones}: quark\_\allowbreak{}down\par \textit{Color}: R\par \textit{Perfiles conjuntos}: 1 & (80,\allowbreak{}54,\allowbreak{}6)\allowbreak{}x1 & (80,\allowbreak{}54,\allowbreak{}6)\allowbreak{}x1 \\
% VI-CATALOG-ROW familias 31
31\par 30|\allowbreak{}-6|\allowbreak{}0|\allowbreak{}0|\allowbreak{}-8|\allowbreak{}+0+-\kern0pt{}-\kern0pt{}-+0+-\kern0pt{}-\kern0pt{}- & \textit{Número de celdas}: 1\par \textit{Celdas (identificadores)}: 54\par \textit{Multisecciones}: quark\_\allowbreak{}down\_\allowbreak{}G4\par \textit{Colecciones}: quark\_\allowbreak{}down\par \textit{Color}: R\par \textit{Perfiles conjuntos}: 1 & (0,\allowbreak{}24,\allowbreak{}0)\allowbreak{}x1 & (80,\allowbreak{}54,\allowbreak{}6)\allowbreak{}x1 \\
% VI-CATALOG-ROW familias 32
32\par 30|\allowbreak{}-6|\allowbreak{}4|\allowbreak{}0|\allowbreak{}-12|\allowbreak{}+++-\kern0pt{}-\kern0pt{}-+++-\kern0pt{}-\kern0pt{}- & \textit{Número de celdas}: 1\par \textit{Celdas (identificadores)}: 65\par \textit{Multisecciones}: quark\_\allowbreak{}up\_\allowbreak{}G3\par \textit{Colecciones}: quark\_\allowbreak{}up\par \textit{Color}: R\par \textit{Perfiles conjuntos}: 1 & (80,\allowbreak{}54,\allowbreak{}6)\allowbreak{}x1 & (80,\allowbreak{}54,\allowbreak{}6)\allowbreak{}x1 \\
% VI-CATALOG-ROW familias 33
33\par 30|\allowbreak{}-6|\allowbreak{}8|\allowbreak{}0|\allowbreak{}-12|\allowbreak{}+++-\kern0pt{}-\kern0pt{}-+++-\kern0pt{}-\kern0pt{}- & \textit{Número de celdas}: 1\par \textit{Celdas (identificadores)}: 44\par \textit{Multisecciones}: neutrino\_\allowbreak{}G3\par \textit{Colecciones}: neutrino\par \textit{Color}: none\par \textit{Perfiles conjuntos}: 1 & (60,\allowbreak{}84,\allowbreak{}28)\allowbreak{}x1 & (80,\allowbreak{}54,\allowbreak{}6)\allowbreak{}x1 \\
% VI-CATALOG-ROW familias 34
34\par 30|\allowbreak{}6|\allowbreak{}-4|\allowbreak{}0|\allowbreak{}-12|\allowbreak{}+++-\kern0pt{}-\kern0pt{}-+++-\kern0pt{}-\kern0pt{}- & \textit{Número de celdas}: 1\par \textit{Celdas (identificadores)}: 68\par \textit{Multisecciones}: quark\_\allowbreak{}down\_\allowbreak{}G3\par \textit{Colecciones}: quark\_\allowbreak{}down\par \textit{Color}: R\par \textit{Perfiles conjuntos}: 1 & (60,\allowbreak{}84,\allowbreak{}28)\allowbreak{}x1 & (80,\allowbreak{}54,\allowbreak{}6)\allowbreak{}x1 \\
% VI-CATALOG-ROW familias 35
35\par 30|\allowbreak{}6|\allowbreak{}-8|\allowbreak{}0|\allowbreak{}-12|\allowbreak{}+++-\kern0pt{}-\kern0pt{}-+++-\kern0pt{}-\kern0pt{}- & \textit{Número de celdas}: 1\par \textit{Celdas (identificadores)}: 17\par \textit{Multisecciones}: quark\_\allowbreak{}up\_\allowbreak{}G3\par \textit{Colecciones}: quark\_\allowbreak{}up\par \textit{Color}: R\par \textit{Perfiles conjuntos}: 1 & (80,\allowbreak{}54,\allowbreak{}6)\allowbreak{}x1 & (80,\allowbreak{}54,\allowbreak{}6)\allowbreak{}x1 \\
% VI-CATALOG-ROW familias 36
36\par 30|\allowbreak{}6|\allowbreak{}0|\allowbreak{}0|\allowbreak{}-12|\allowbreak{}+++-\kern0pt{}-\kern0pt{}-+++-\kern0pt{}-\kern0pt{}- & \textit{Número de celdas}: 2\par \textit{Celdas (identificadores)}: 38,\allowbreak{}58\par \textit{Multisecciones}: neutrino\_\allowbreak{}G2 | neutrino\_\allowbreak{}G3\par \textit{Colecciones}: neutrino\par \textit{Color}: none\par \textit{Perfiles conjuntos}: 2 & (60,\allowbreak{}78,\allowbreak{}25)\allowbreak{}x1; (80,\allowbreak{}54,\allowbreak{}6)\allowbreak{}x1 & (80,\allowbreak{}54,\allowbreak{}6)\allowbreak{}x2 \\
% VI-CATALOG-ROW familias 37
37\par 30|\allowbreak{}6|\allowbreak{}0|\allowbreak{}0|\allowbreak{}-8|\allowbreak{}+++-0-+++-0- & \textit{Número de celdas}: 1\par \textit{Celdas (identificadores)}: 78\par \textit{Multisecciones}: neutrino\_\allowbreak{}G4\par \textit{Colecciones}: neutrino\par \textit{Color}: none\par \textit{Perfiles conjuntos}: 1 & (0,\allowbreak{}24,\allowbreak{}0)\allowbreak{}x1 & (80,\allowbreak{}54,\allowbreak{}6)\allowbreak{}x1 \\
% VI-CATALOG-ROW familias 38
38\par 30|\allowbreak{}6|\allowbreak{}0|\allowbreak{}2|\allowbreak{}-8|\allowbreak{}+++-\kern0pt{}-0+++-\kern0pt{}-0 & \textit{Número de celdas}: 1\par \textit{Celdas (identificadores)}: 27\par \textit{Multisecciones}: charged\_\allowbreak{}lepton\_\allowbreak{}G4\par \textit{Colecciones}: charged\_\allowbreak{}lepton\par \textit{Color}: none\par \textit{Perfiles conjuntos}: 1 & (40,\allowbreak{}54,\allowbreak{}6)\allowbreak{}x1 & (56,\allowbreak{}54,\allowbreak{}5)\allowbreak{}x1 \\
% VI-CATALOG-ROW familias 39
39\par 30|\allowbreak{}6|\allowbreak{}0|\allowbreak{}2|\allowbreak{}-8|\allowbreak{}+++0-\kern0pt{}-+++0-\kern0pt{}- & \textit{Número de celdas}: 1\par \textit{Celdas (identificadores)}: 48\par \textit{Multisecciones}: quark\_\allowbreak{}down\_\allowbreak{}G2\par \textit{Colecciones}: quark\_\allowbreak{}down\par \textit{Color}: R\par \textit{Perfiles conjuntos}: 1 & (0,\allowbreak{}24,\allowbreak{}0)\allowbreak{}x1 & (4,\allowbreak{}54,\allowbreak{}2)\allowbreak{}x1 \\
% VI-CATALOG-ROW familias 40
40\par 30|\allowbreak{}6|\allowbreak{}4|\allowbreak{}0|\allowbreak{}-12|\allowbreak{}+++-\kern0pt{}-\kern0pt{}-+++-\kern0pt{}-\kern0pt{}- & \textit{Número de celdas}: 1\par \textit{Celdas (identificadores)}: 7\par \textit{Multisecciones}: charged\_\allowbreak{}lepton\_\allowbreak{}G4\par \textit{Colecciones}: charged\_\allowbreak{}lepton\par \textit{Color}: none\par \textit{Perfiles conjuntos}: 1 & (60,\allowbreak{}78,\allowbreak{}25)\allowbreak{}x1 & (80,\allowbreak{}54,\allowbreak{}6)\allowbreak{}x1 \\
% VI-CATALOG-ROW familias 41
41\par 30|\allowbreak{}6|\allowbreak{}8|\allowbreak{}0|\allowbreak{}-12|\allowbreak{}+++-\kern0pt{}-\kern0pt{}-+++-\kern0pt{}-\kern0pt{}- & \textit{Número de celdas}: 1\par \textit{Celdas (identificadores)}: 28\par \textit{Multisecciones}: quark\_\allowbreak{}down\_\allowbreak{}G4\par \textit{Colecciones}: quark\_\allowbreak{}down\par \textit{Color}: R\par \textit{Perfiles conjuntos}: 1 & (80,\allowbreak{}54,\allowbreak{}7)\allowbreak{}x1 & (80,\allowbreak{}54,\allowbreak{}6)\allowbreak{}x1 \\
% VI-CATALOG-ROW familias 42
42\par 34|\allowbreak{}0|\allowbreak{}-12|\allowbreak{}-2|\allowbreak{}-8|\allowbreak{}0++-\kern0pt{}-\kern0pt{}-0++-\kern0pt{}-\kern0pt{}- & \textit{Número de celdas}: 1\par \textit{Celdas (identificadores)}: 45\par \textit{Multisecciones}: charged\_\allowbreak{}lepton\_\allowbreak{}G4\par \textit{Colecciones}: charged\_\allowbreak{}lepton\par \textit{Color}: none\par \textit{Perfiles conjuntos}: 1 & (80,\allowbreak{}66,\allowbreak{}4)\allowbreak{}x1 & (56,\allowbreak{}54,\allowbreak{}5)\allowbreak{}x1 \\
% VI-CATALOG-ROW familias 43
43\par 34|\allowbreak{}0|\allowbreak{}-12|\allowbreak{}0|\allowbreak{}-12|\allowbreak{}+++-\kern0pt{}-\kern0pt{}-+++-\kern0pt{}-\kern0pt{}- & \textit{Número de celdas}: 2\par \textit{Celdas (identificadores)}: 57,\allowbreak{}76\par \textit{Multisecciones}: charged\_\allowbreak{}lepton\_\allowbreak{}G2 | neutrino\_\allowbreak{}G4\par \textit{Colecciones}: charged\_\allowbreak{}lepton | neutrino\par \textit{Color}: none\par \textit{Perfiles conjuntos}: 2 & (40,\allowbreak{}84,\allowbreak{}30)\allowbreak{}x1; (80,\allowbreak{}66,\allowbreak{}2)\allowbreak{}x1 & (80,\allowbreak{}54,\allowbreak{}6)\allowbreak{}x2 \\
% VI-CATALOG-ROW familias 44
44\par 34|\allowbreak{}0|\allowbreak{}-16|\allowbreak{}0|\allowbreak{}-8|\allowbreak{}+++-0-+++-0- & \textit{Número de celdas}: 1\par \textit{Celdas (identificadores)}: 37\par \textit{Multisecciones}: charged\_\allowbreak{}lepton\_\allowbreak{}G4\par \textit{Colecciones}: charged\_\allowbreak{}lepton\par \textit{Color}: none\par \textit{Perfiles conjuntos}: 1 & (40,\allowbreak{}24,\allowbreak{}2)\allowbreak{}x1 & (80,\allowbreak{}54,\allowbreak{}6)\allowbreak{}x1 \\
% VI-CATALOG-ROW familias 45
45\par 34|\allowbreak{}0|\allowbreak{}-16|\allowbreak{}0|\allowbreak{}-8|\allowbreak{}+0+-\kern0pt{}-\kern0pt{}-+0+-\kern0pt{}-\kern0pt{}- & \textit{Número de celdas}: 1\par \textit{Celdas (identificadores)}: 15\par \textit{Multisecciones}: quark\_\allowbreak{}up\_\allowbreak{}G3\par \textit{Colecciones}: quark\_\allowbreak{}up\par \textit{Color}: B\par \textit{Perfiles conjuntos}: 1 & (40,\allowbreak{}78,\allowbreak{}27)\allowbreak{}x1 & (80,\allowbreak{}54,\allowbreak{}6)\allowbreak{}x1 \\
% VI-CATALOG-ROW familias 46
46\par 34|\allowbreak{}0|\allowbreak{}-20|\allowbreak{}-2|\allowbreak{}-8|\allowbreak{}0++-\kern0pt{}-\kern0pt{}-0++-\kern0pt{}-\kern0pt{}- & \textit{Número de celdas}: 1\par \textit{Celdas (identificadores)}: 35\par \textit{Multisecciones}: quark\_\allowbreak{}up\_\allowbreak{}G3\par \textit{Colecciones}: quark\_\allowbreak{}up\par \textit{Color}: B\par \textit{Perfiles conjuntos}: 1 & (100,\allowbreak{}66,\allowbreak{}0)\allowbreak{}x1 & (56,\allowbreak{}54,\allowbreak{}5)\allowbreak{}x1 \\
% VI-CATALOG-ROW familias 47
47\par 34|\allowbreak{}0|\allowbreak{}-4|\allowbreak{}-2|\allowbreak{}-8|\allowbreak{}++0-\kern0pt{}-\kern0pt{}-++0-\kern0pt{}-\kern0pt{}- & \textit{Número de celdas}: 1\par \textit{Celdas (identificadores)}: 56\par \textit{Multisecciones}: neutrino\_\allowbreak{}G3\par \textit{Colecciones}: neutrino\par \textit{Color}: none\par \textit{Perfiles conjuntos}: 1 & (40,\allowbreak{}24,\allowbreak{}2)\allowbreak{}x1 & (36,\allowbreak{}54,\allowbreak{}4)\allowbreak{}x1 \\
% VI-CATALOG-ROW familias 48
48\par 34|\allowbreak{}0|\allowbreak{}-4|\allowbreak{}0|\allowbreak{}-12|\allowbreak{}+++-\kern0pt{}-\kern0pt{}-+++-\kern0pt{}-\kern0pt{}- & \textit{Número de celdas}: 3\par \textit{Celdas (identificadores)}: 6,\allowbreak{}25,\allowbreak{}46\par \textit{Multisecciones}: charged\_\allowbreak{}lepton\_\allowbreak{}G2 | neutrino\_\allowbreak{}G4 | quark\_\allowbreak{}down\_\allowbreak{}G4\par \textit{Colecciones}: charged\_\allowbreak{}lepton | neutrino | quark\_\allowbreak{}down\par \textit{Color}: B | none\par \textit{Perfiles conjuntos}: 2 & (80,\allowbreak{}66,\allowbreak{}2)\allowbreak{}x1; (80,\allowbreak{}66,\allowbreak{}3)\allowbreak{}x2 & (80,\allowbreak{}54,\allowbreak{}6)\allowbreak{}x3 \\
% VI-CATALOG-ROW familias 49
49\par 34|\allowbreak{}0|\allowbreak{}-4|\allowbreak{}2|\allowbreak{}-8|\allowbreak{}+++-\kern0pt{}-0+++-\kern0pt{}-0 & \textit{Número de celdas}: 1\par \textit{Celdas (identificadores)}: 67\par \textit{Multisecciones}: quark\_\allowbreak{}up\_\allowbreak{}G3\par \textit{Colecciones}: quark\_\allowbreak{}up\par \textit{Color}: G\par \textit{Perfiles conjuntos}: 1 & (100,\allowbreak{}66,\allowbreak{}2)\allowbreak{}x1 & (56,\allowbreak{}54,\allowbreak{}5)\allowbreak{}x1 \\
% VI-CATALOG-ROW familias 50
50\par 34|\allowbreak{}0|\allowbreak{}-8|\allowbreak{}-2|\allowbreak{}-8|\allowbreak{}++0-\kern0pt{}-\kern0pt{}-++0-\kern0pt{}-\kern0pt{}- & \textit{Número de celdas}: 1\par \textit{Celdas (identificadores)}: 66\par \textit{Multisecciones}: quark\_\allowbreak{}down\_\allowbreak{}G3\par \textit{Colecciones}: quark\_\allowbreak{}down\par \textit{Color}: B\par \textit{Perfiles conjuntos}: 1 & (40,\allowbreak{}78,\allowbreak{}27)\allowbreak{}x1 & (80,\allowbreak{}54,\allowbreak{}6)\allowbreak{}x1 \\
% VI-CATALOG-ROW familias 51
51\par 34|\allowbreak{}0|\allowbreak{}-8|\allowbreak{}-2|\allowbreak{}-8|\allowbreak{}0++-\kern0pt{}-\kern0pt{}-0++-\kern0pt{}-\kern0pt{}- & \textit{Número de celdas}: 1\par \textit{Celdas (identificadores)}: 77\par \textit{Multisecciones}: charged\_\allowbreak{}lepton\_\allowbreak{}G4\par \textit{Colecciones}: charged\_\allowbreak{}lepton\par \textit{Color}: none\par \textit{Perfiles conjuntos}: 1 & (80,\allowbreak{}66,\allowbreak{}3)\allowbreak{}x1 & (56,\allowbreak{}54,\allowbreak{}5)\allowbreak{}x1 \\
% VI-CATALOG-ROW familias 52
52\par 34|\allowbreak{}0|\allowbreak{}-8|\allowbreak{}0|\allowbreak{}-12|\allowbreak{}+++-\kern0pt{}-\kern0pt{}-+++-\kern0pt{}-\kern0pt{}- & \textit{Número de celdas}: 1\par \textit{Celdas (identificadores)}: 36\par \textit{Multisecciones}: quark\_\allowbreak{}down\_\allowbreak{}G4\par \textit{Colecciones}: quark\_\allowbreak{}down\par \textit{Color}: G\par \textit{Perfiles conjuntos}: 1 & (40,\allowbreak{}84,\allowbreak{}30)\allowbreak{}x1 & (80,\allowbreak{}54,\allowbreak{}6)\allowbreak{}x1 \\
% VI-CATALOG-ROW familias 53
53\par 34|\allowbreak{}0|\allowbreak{}-8|\allowbreak{}0|\allowbreak{}-8|\allowbreak{}+0+-\kern0pt{}-\kern0pt{}-+0+-\kern0pt{}-\kern0pt{}- & \textit{Número de celdas}: 1\par \textit{Celdas (identificadores)}: 5\par \textit{Multisecciones}: charged\_\allowbreak{}lepton\_\allowbreak{}G4\par \textit{Colecciones}: charged\_\allowbreak{}lepton\par \textit{Color}: none\par \textit{Perfiles conjuntos}: 1 & (40,\allowbreak{}24,\allowbreak{}2)\allowbreak{}x1 & (80,\allowbreak{}54,\allowbreak{}6)\allowbreak{}x1 \\
% VI-CATALOG-ROW familias 54
54\par 34|\allowbreak{}0|\allowbreak{}0|\allowbreak{}-2|\allowbreak{}-8|\allowbreak{}++0-\kern0pt{}-\kern0pt{}-++0-\kern0pt{}-\kern0pt{}- & \textit{Número de celdas}: 1\par \textit{Celdas (identificadores)}: 26\par \textit{Multisecciones}: neutrino\_\allowbreak{}G3\par \textit{Colecciones}: neutrino\par \textit{Color}: none\par \textit{Perfiles conjuntos}: 1 & (40,\allowbreak{}78,\allowbreak{}27)\allowbreak{}x1 & (80,\allowbreak{}54,\allowbreak{}6)\allowbreak{}x1 \\
% VI-CATALOG-ROW familias 55
55\par 34|\allowbreak{}0|\allowbreak{}0|\allowbreak{}0|\allowbreak{}-8|\allowbreak{}+0+-\kern0pt{}-\kern0pt{}-+0+-\kern0pt{}-\kern0pt{}- & \textit{Número de celdas}: 1\par \textit{Celdas (identificadores)}: 47\par \textit{Multisecciones}: quark\_\allowbreak{}up\_\allowbreak{}G3\par \textit{Colecciones}: quark\_\allowbreak{}up\par \textit{Color}: G\par \textit{Perfiles conjuntos}: 1 & (40,\allowbreak{}78,\allowbreak{}27)\allowbreak{}x1 & (80,\allowbreak{}54,\allowbreak{}6)\allowbreak{}x1 \\
% VI-CATALOG-ROW familias 56
56\par 34|\allowbreak{}0|\allowbreak{}0|\allowbreak{}2|\allowbreak{}-8|\allowbreak{}+++0-\kern0pt{}-+++0-\kern0pt{}- & \textit{Número de celdas}: 1\par \textit{Celdas (identificadores)}: 16\par \textit{Multisecciones}: quark\_\allowbreak{}down\_\allowbreak{}G3\par \textit{Colecciones}: quark\_\allowbreak{}down\par \textit{Color}: G\par \textit{Perfiles conjuntos}: 1 & (40,\allowbreak{}24,\allowbreak{}2)\allowbreak{}x1 & (36,\allowbreak{}54,\allowbreak{}4)\allowbreak{}x1 \\
\end{longtable}
\endgroup
\subsection{Clasificación de los 23 sectores documentados}

Estas filas reúnen sectores explícitos, residencias y representación. La clave
de estado remite a la leyenda contigua y conserva el estatuto de la fuente.
En particular, un operador interno exacto, una evaluación condicionada y una
realización escalar cerrada se mantienen como afirmaciones distintas. La tabla
no actualiza su estatuto ni transforma un registro condicionado en una predicción.

\begingroup
\setlength{\tabcolsep}{4pt}
\setlength{\extrarowheight}{4pt}
\renewcommand{\arraystretch}{1.13}
\begin{longtable}{@{}>{\raggedright\arraybackslash}p{\dimexpr 0.08860759\linewidth-2\tabcolsep\relax}>{\raggedright\arraybackslash}p{\dimexpr 0.41772152\linewidth-2\tabcolsep\relax}>{\raggedright\arraybackslash}p{\dimexpr 0.49367089\linewidth-2\tabcolsep\relax}@{}}
\caption{Estados originales de las clases.}\label{tab:vi-estados-clases}\\
\toprule
Clave & Estado literal & Lectura documental \\
\midrule
\endfirsthead
\multicolumn{3}{l}{\textit{Continuación}}\\
\toprule
Clave & Estado literal & Lectura documental \\
\midrule
\endhead
\midrule
\multicolumn{3}{r}{Continúa en la página siguiente}\\
\endfoot
\bottomrule
\endlastfoot
C01 & CLOSED\_\allowbreak{}ELECTRON\_\allowbreak{}TWO\_\allowbreak{}WAY & Electrón: cierre bidireccional documentado. \\
C02 & CLOSED\_\allowbreak{}NULL\_\allowbreak{}CHART & Coordenatización nula cerrada en este corte. \\
C03 & NOT\_\allowbreak{}A\_\allowbreak{}UNIVERSAL\_\allowbreak{}MASS\_\allowbreak{}SCALAR & La clase no se expresa como un escalar universal de masa. \\
C04 & OPEN\_\allowbreak{}FLAVOR\_\allowbreak{}COLOR\_\allowbreak{}SCALARIZATION & Escalarización de sabor y color abierta en este corte. \\
C05 & OPEN\_\allowbreak{}MIXING\_\allowbreak{}AND\_\allowbreak{}SCALE & Mezcla y escala abiertas en este corte. \\
C06 & OPEN\_\allowbreak{}NAMED\_\allowbreak{}STATE\_\allowbreak{}SCALARIZATION & Escalarización del estado individualizado abierta en este corte. \\
C07 & OPEN\_\allowbreak{}PHYSICAL\_\allowbreak{}SCALARIZATION & Escalarización física abierta en este corte. \\
C08 & OPEN\_\allowbreak{}ROUTE\_\allowbreak{}AND\_\allowbreak{}SCALARIZATION & Ruta y escalarización abiertas en este corte. \\
\end{longtable}
\endgroup

\begingroup
\setlength{\tabcolsep}{4pt}
\setlength{\extrarowheight}{4pt}
\renewcommand{\arraystretch}{1.08}
\begin{longtable}{@{}>{\raggedright\arraybackslash}p{\dimexpr 0.17307692\linewidth-2\tabcolsep\relax}>{\raggedright\arraybackslash}p{\dimexpr 0.38461538\linewidth-2\tabcolsep\relax}>{\raggedright\arraybackslash}p{\dimexpr 0.33333333\linewidth-2\tabcolsep\relax}>{\raggedright\arraybackslash}p{\dimexpr 0.10897436\linewidth-2\tabcolsep\relax}@{}}
\caption{Las 23 clases del catálogo, en su orden documental.}\label{tab:vi-clases}\\
\toprule
Clase & Representantes y residencia & Incidencia y representación & Estado \\
\midrule
\endfirsthead
\multicolumn{4}{l}{\textit{Continuación}}\\
\toprule
Clase & Representantes y residencia & Incidencia y representación & Estado \\
\midrule
\endhead
\midrule
\multicolumn{4}{r}{Continúa en la página siguiente}\\
\endfoot
\bottomrule
\endlastfoot
% VI-CATALOG-ROW clases 1
SM-LC-1\par lepton\_\allowbreak{}cargado & \textit{Representantes}: e− / e+\par \textit{Residencia}: centro (5,\allowbreak{}5) y antisección & \textit{Celdas seleccionadas}: 1\par \textit{Multisecciones}: charged\_\allowbreak{}lepton\_\allowbreak{}G0\par \textit{Carga}: ±1 e\par \textit{Representación}: sin color & C01 \\
% VI-CATALOG-ROW clases 2
SM-LC-2\par lepton\_\allowbreak{}cargado & \textit{Representantes}: μ− / μ+\par \textit{Residencia}: multisección leptón cargado & \textit{Celdas seleccionadas}: 8\par \textit{Multisecciones}: charged\_\allowbreak{}lepton\_\allowbreak{}G2\par \textit{Carga}: ±1 e\par \textit{Representación}: sin color & C07 \\
% VI-CATALOG-ROW clases 3
SM-LC-3\par lepton\_\allowbreak{}cargado & \textit{Representantes}: τ− / τ+\par \textit{Residencia}: multisección leptón cargado & \textit{Celdas seleccionadas}: 16\par \textit{Multisecciones}: charged\_\allowbreak{}lepton\_\allowbreak{}G4\par \textit{Carga}: ±1 e\par \textit{Representación}: sin color & C07 \\
% VI-CATALOG-ROW clases 4
SM-NU-1\par neutrino & \textit{Representantes}: ν\_\allowbreak{}e / anti-ν\_\allowbreak{}e\par \textit{Residencia}: multisecciones neutrino & \textit{Celdas seleccionadas}: 20\par \textit{Multisecciones}: neutrino\_\allowbreak{}G1 | neutrino\_\allowbreak{}G2 | neutrino\_\allowbreak{}G3 | neutrino\_\allowbreak{}G4\par \textit{Carga}: 0\par \textit{Representación}: sin color & C05 \\
% VI-CATALOG-ROW clases 5
SM-NU-2\par neutrino & \textit{Representantes}: ν\_\allowbreak{}μ / anti-ν\_\allowbreak{}μ\par \textit{Residencia}: multisecciones neutrino & \textit{Celdas seleccionadas}: 20\par \textit{Multisecciones}: neutrino\_\allowbreak{}G1 | neutrino\_\allowbreak{}G2 | neutrino\_\allowbreak{}G3 | neutrino\_\allowbreak{}G4\par \textit{Carga}: 0\par \textit{Representación}: sin color & C05 \\
% VI-CATALOG-ROW clases 6
SM-NU-3\par neutrino & \textit{Representantes}: ν\_\allowbreak{}τ / anti-ν\_\allowbreak{}τ\par \textit{Residencia}: multisecciones neutrino & \textit{Celdas seleccionadas}: 20\par \textit{Multisecciones}: neutrino\_\allowbreak{}G1 | neutrino\_\allowbreak{}G2 | neutrino\_\allowbreak{}G3 | neutrino\_\allowbreak{}G4\par \textit{Carga}: 0\par \textit{Representación}: sin color & C05 \\
% VI-CATALOG-ROW clases 7
SM-Q-U1\par quark & \textit{Representantes}: u / anti-u\par \textit{Residencia}: multisección tipo arriba & \textit{Celdas seleccionadas}: 16\par \textit{Multisecciones}: quark\_\allowbreak{}up\_\allowbreak{}G1 | quark\_\allowbreak{}up\_\allowbreak{}G3\par \textit{Carga}: ±2/\allowbreak{}3 e\par \textit{Representación}: R,\allowbreak{}G,\allowbreak{}B = órbita Z/\allowbreak{}3 & C04 \\
% VI-CATALOG-ROW clases 8
SM-Q-D1\par quark & \textit{Representantes}: d / anti-d\par \textit{Residencia}: multisección tipo abajo & \textit{Celdas seleccionadas}: 20\par \textit{Multisecciones}: quark\_\allowbreak{}down\_\allowbreak{}G1 | quark\_\allowbreak{}down\_\allowbreak{}G2 | quark\_\allowbreak{}down\_\allowbreak{}G3 | quark\_\allowbreak{}down\_\allowbreak{}G4\par \textit{Carga}: ±1/\allowbreak{}3 e\par \textit{Representación}: R,\allowbreak{}G,\allowbreak{}B = órbita Z/\allowbreak{}3 & C04 \\
% VI-CATALOG-ROW clases 9
SM-Q-U2\par quark & \textit{Representantes}: c / anti-c\par \textit{Residencia}: multisección tipo arriba & \textit{Celdas seleccionadas}: 16\par \textit{Multisecciones}: quark\_\allowbreak{}up\_\allowbreak{}G1 | quark\_\allowbreak{}up\_\allowbreak{}G3\par \textit{Carga}: ±2/\allowbreak{}3 e\par \textit{Representación}: R,\allowbreak{}G,\allowbreak{}B = órbita Z/\allowbreak{}3 & C04 \\
% VI-CATALOG-ROW clases 10
SM-Q-D2\par quark & \textit{Representantes}: s / anti-s\par \textit{Residencia}: multisección tipo abajo & \textit{Celdas seleccionadas}: 20\par \textit{Multisecciones}: quark\_\allowbreak{}down\_\allowbreak{}G1 | quark\_\allowbreak{}down\_\allowbreak{}G2 | quark\_\allowbreak{}down\_\allowbreak{}G3 | quark\_\allowbreak{}down\_\allowbreak{}G4\par \textit{Carga}: ±1/\allowbreak{}3 e\par \textit{Representación}: R,\allowbreak{}G,\allowbreak{}B = órbita Z/\allowbreak{}3 & C04 \\
% VI-CATALOG-ROW clases 11
SM-Q-U3\par quark & \textit{Representantes}: t / anti-t\par \textit{Residencia}: multisección tipo arriba & \textit{Celdas seleccionadas}: 16\par \textit{Multisecciones}: quark\_\allowbreak{}up\_\allowbreak{}G1 | quark\_\allowbreak{}up\_\allowbreak{}G3\par \textit{Carga}: ±2/\allowbreak{}3 e\par \textit{Representación}: R,\allowbreak{}G,\allowbreak{}B = órbita Z/\allowbreak{}3 & C04 \\
% VI-CATALOG-ROW clases 12
SM-Q-D3\par quark & \textit{Representantes}: b / anti-b\par \textit{Residencia}: multisección tipo abajo & \textit{Celdas seleccionadas}: 20\par \textit{Multisecciones}: quark\_\allowbreak{}down\_\allowbreak{}G1 | quark\_\allowbreak{}down\_\allowbreak{}G2 | quark\_\allowbreak{}down\_\allowbreak{}G3 | quark\_\allowbreak{}down\_\allowbreak{}G4\par \textit{Carga}: ±1/\allowbreak{}3 e\par \textit{Representación}: R,\allowbreak{}G,\allowbreak{}B = órbita Z/\allowbreak{}3 & C04 \\
% VI-CATALOG-ROW clases 13
SM-GA-EM\par boson\_\allowbreak{}gauge & \textit{Representantes}: γ\par \textit{Residencia}: sector gauge nulo & \textit{Celdas seleccionadas}: 0\par \textit{Multisecciones}: \textemdash{}\par \textit{Carga}: 0\par \textit{Representación}: sin color & C02 \\
% VI-CATALOG-ROW clases 14
SM-GA-GL\par boson\_\allowbreak{}gauge & \textit{Representantes}: g\par \textit{Residencia}: sector gauge de color & \textit{Celdas seleccionadas}: 0\par \textit{Multisecciones}: \textemdash{}\par \textit{Carga}: 0\par \textit{Representación}: octete gauge & C02 \\
% VI-CATALOG-ROW clases 15
SM-GA-W\par boson\_\allowbreak{}gauge & \textit{Representantes}: W+ / W−\par \textit{Residencia}: sector gauge masivo & \textit{Celdas seleccionadas}: 0\par \textit{Multisecciones}: \textemdash{}\par \textit{Carga}: ±1 e\par \textit{Representación}: sin color & C08 \\
% VI-CATALOG-ROW clases 16
SM-GA-Z\par boson\_\allowbreak{}gauge & \textit{Representantes}: Z0\par \textit{Residencia}: sector gauge masivo & \textit{Celdas seleccionadas}: 0\par \textit{Multisecciones}: \textemdash{}\par \textit{Carga}: 0\par \textit{Representación}: sin color & C08 \\
% VI-CATALOG-ROW clases 17
SM-H-1\par escalar & \textit{Representantes}: H\par \textit{Residencia}: sector escalar & \textit{Celdas seleccionadas}: 0\par \textit{Multisecciones}: \textemdash{}\par \textit{Carga}: 0\par \textit{Representación}: sin color & C08 \\
% VI-CATALOG-ROW clases 18
COMP-MESON\par meson & \textit{Representantes}: π0, π±, K0, K±, η, ρ, φ, J/\allowbreak{}ψ, Υ…\par \textit{Residencia}: producto de sección y antisección & \textit{Celdas seleccionadas}: 0\par \textit{Multisecciones}: \textemdash{}\par \textit{Carga}: dependiente del estado compuesto\par \textit{Representación}: singlete u otras representaciones & C06 \\
% VI-CATALOG-ROW clases 19
COMP-BARYON\par barion & \textit{Representantes}: p, n y configuraciones bariónicas\par \textit{Residencia}: producto de tres secciones & \textit{Celdas seleccionadas}: 0\par \textit{Multisecciones}: \textemdash{}\par \textit{Carga}: dependiente del estado compuesto\par \textit{Representación}: singlete físico & C06 \\
% VI-CATALOG-ROW clases 20
TOPO-FIB\par anyon\_\allowbreak{}no\_\allowbreak{}abeliano & \textit{Representantes}: 1, τ\par \textit{Residencia}: sector topológico & \textit{Celdas seleccionadas}: 0\par \textit{Multisecciones}: \textemdash{}\par \textit{Carga}: definida por el codominio topológico\par \textit{Representación}: no aplica & C03 \\
% VI-CATALOG-ROW clases 21
TOPO-ISING\par sector\_\allowbreak{}majorana & \textit{Representantes}: 1, ψ, σ\par \textit{Residencia}: sector topológico / cero & \textit{Celdas seleccionadas}: 0\par \textit{Multisecciones}: \textemdash{}\par \textit{Carga}: dependiente de la realización\par \textit{Representación}: sin color & C03 \\
% VI-CATALOG-ROW clases 22
TOPO-Z6\par anyones\_\allowbreak{}abelianos & \textit{Representantes}: bosón, fermión y cuatro sectores trenzados\par \textit{Residencia}: sector topológico Z/\allowbreak{}6 & \textit{Celdas seleccionadas}: 0\par \textit{Multisecciones}: \textemdash{}\par \textit{Carga}: Z/\allowbreak{}6 definida por el trenzado\par \textit{Representación}: sin color & C03 \\
% VI-CATALOG-ROW clases 23
VIRT-1\par virtual & \textit{Representantes}: secciones virtuales\par \textit{Residencia}: sistema inverso TPK & \textit{Celdas seleccionadas}: 0\par \textit{Multisecciones}: \textemdash{}\par \textit{Carga}: dependiente del sector de prolongación\par \textit{Representación}: según sector & C03 \\
\end{longtable}
\endgroup

\input{sections/12_catalogo_datos_companero.tex}
\end{document}
